% !TEX root = ../notes.tex

\documentclass[../notes.tex]{subfiles}

\begin{document}

\section{March 30}
This talk was given by Zhiwei Yun.

\subsection{Review}
Let's recall what was covered. Fix a reductive group $G$ over an algebraically closed field, and let $W$ be the Weyl group. For the past few weeks, we have been discussing the strata of a reductive group $G$, which today we will denote $\op{St}(G)$. It was related to both conjugacy classes in the Weyl group and the irreducible representations of the Weyl group $W$. In particular, we had the following two maps.
\begin{itemize}
	\item Taking some inductions produces an embedding map
	\[\op{Str}(G)\to\op{Irr}(G).\]
	The image is known as the $2$-special representations and is denoted $\op{Irr}_*W$. This map admits a section given by the Springer correspondence, which maps to $\op{Conj}(U_p)$ (where $p$ is the characteristic), which then embeds into $\op{Irr}_*W$.
	\item There was a complicated map $\op{cl}(W)\to\op{Irr}(G)$. We will give a construction in characteristic zero shortly.
	\item The complicated map $\op{cl}(W)\to\op{Irr}(G)$ admits two canonical sections, given by choosing either the most or least elliptic elements of $G$.
\end{itemize}
Let's specialize to characteristic zero. Then let $\mc N$ be the nilpotent orbits of the Lie algebra $\mf g$, and we are able to gives maps to and from $\mc N$ as above.
\begin{itemize}
	\item By the Springer correspondence, there is a map $\mc N\to\op{Irr}(W)$ given by sending a nilpotent orbit $\OO$ to $\mathrm H^{\mathrm{top}}(\mc B_e)$. Conversely, the Springer correspondence associates an orbit and a local system to each irreducible representation in $\op{Irr}W$, so one can forget the local system to recover a nilpotent orbit.
	\item There is also a surjection $\op{cl}(W)\to\mc N$. Namely, given a class in $W$, let $w$ be a representative of minimal length. Then it is a highly nontrivial theorem of Lusztig that $BwB\cap\mc U$ admits intersection with a unique minimal nilpotent orbit. (Here, $\mc U$ is the unipotent variety.) There are once again sections $\mc N\to\op{cl}(W)$ given by choosing the most or least elliptic elements.
\end{itemize}
Today, we will focus on the maps to and from $\op{cl}(W)$.

\subsection{The Flag Variety}
For the rest of today, we fix a reductive group $G$ over $\CC$. Here is some notation for today.
\begin{definition}[loop group]
	Fix a reductive group $G$ over $\CC$. Then the \textit{loop group} $LG$ is $LG(\CC)\coloneqq G(\CC((t)))$. We may also write $L^+G\coloneqq G(\CC[[t]])$, and then there is a Grassmannian $\mathrm{Gr}\coloneqq LG/L^+G$.
\end{definition}
\begin{remark}
	The group $LG$ is an infinite-dimensional Lie group. Notably, it may admit different connected components if $G$ is not simply connected.
\end{remark}
\begin{definition}[flag variety]
	Fix a reductive group $G$ over $\CC$. Then the \textit{Iwahori subgroup} $I$ consists of those $g\in L^+G$ such that specializing $t\to0$ places $g$ in $B$. We then define the \textit{flag variety} as
	\[\mathrm{Fl}\coloneqq LG/I.\]
\end{definition}
\begin{remark}
	This flag variety turns out to admit a stratification $\bigcup_{w\in\widetilde W}\op{Fl}_w$, where $\widetilde W\coloneqq X_*(T)\rtimes W$ is the affine Weyl group.
\end{remark}
\begin{notation}
	There are affine Springer fibers defined as follows: for $\gamma\in\mf g((t))$, we define
	\[\mathrm{Fl}_\gamma\coloneqq\{g\in LG/I:\op{Ad}_{g^{-1}}\gamma\in\{X\in\mf g[[t]]:X_{t=0}\in\mf b\}\}.\]
	This is exactly analogous to the usual Springer fiber $\mc B_e=\{g\in G/B:\op{Ad}_{g^{-1}}e\in\mf b\}$. While we're here, we set
	\[\mathrm{Gr}_\gamma\coloneqq\{g\in LG/L^+G:\op{Ad}_{g^{-1}}\gamma\in\mf g[[t]]\}.\]
\end{notation}
\begin{remark}
	It turns out that $\op{Fl}_\gamma$ and $\op{Gr}_\gamma$ are nonempty if and only if $\gamma$ is integral, which means that $\gamma$ can be conjugated into $\mf g[[t]]$. (This has no content for $\op{Gr}_\gamma$.) For $\op{GL}_n$, being integral amounts to saying that the characteristic polynomial lives in $\CC[[t]]$.
\end{remark}
\begin{remark}
	The fiber of the map $\mathrm{Fl}_\gamma\to\mathrm{Gr}_\gamma$ looks like $\mc B_e$, which varies over the base $\mathrm{Gr}_\gamma$.
\end{remark}
We will not be interested in all $\gamma$s, even those which are integral.
\begin{remark}
	If $\gamma$ is regular and semisimple, then $\op{Fl}_\gamma$ and $\op{Gr}_\gamma$ are finite-dimensional, but they may have infinitely many irreducible components. In this case, Kazhdan--Lusztig (modifying an idea of Spaltenstein) showed that the $\op{Fl}_\gamma$s are equidimensional with a conjectural formula for the dimension, which was proven by Bezrukavnikov; Ng\^o showed that the $\op{Gr}_\gamma$s are equidimensional.
\end{remark}
\begin{definition}[topologically nilpotent]
	Fix a Lie algebra $\mf g$ over $\CC$, and choose $\gamma\in\mf g[[t]]$. Then $\gamma$ is \textit{topologically nilpotent} (for the $t$-adic topology) if and only if it can be conjugated into $\{X\in\mf g[[t]]:X_{t=0}\in\mf n\}$.
\end{definition}
\begin{example}
	For $\op{GL}_n$, being topologically nilpotent amounts to saying that the characteristic polynomial has vanishing constant term.
\end{example}

\subsection{Kazhdan--Lusztig Maps}
We are now able to define a map in one direction.
\begin{definition}[Kazhdan--Lusztig map]
	There is a Kazhdan--Lusztig map $\op{KL}\colon\mc N\to\op{cl}(W)$ defined as follows. Given $e\in\mc N$, choose a generic power series
	\[\gamma=e+X_1t+X_2t^2+\cdots\in\mf g[[t]].\]
	(The genericity is rather non-explicit, even once we assume that $\gamma$ is topologically nilpotent---which is automatic here---and regular semisimple.)
	% (For our genericity, $e$ is automatically topologically nilpotent, the genericity is not required, but it is required in the regular semisimple case.)
	Then we define the class $[w]_\gamma\in\op{cl}(W)$ to be given by the splitting type of $\gamma$. For generic $\gamma$, this $[w]_\gamma$ does not depend on the choice of $\gamma$.
\end{definition}
\begin{example}
	For $G=\op{GL}_n$, the characteristic polynomial of $\gamma$ can be written as a product of irreducible polynomials in $\CC((t))[x]$. This factorization gives a partition of $n$, which corresponds to a conjugacy class in $S_n$.
\end{example}
\begin{remark}
	For general $G$, the centralizers $C_{LG}(\gamma)$ are some tori over $\CC((t))$, which is classified by $\mathrm H^1(\CC((t));W)$, but the latter is $\op{cl}(W)$ because $\op{Gal}(\overline{\CC((t))}/\CC((t)))=\widehat{\ZZ}$. This explains how to recover a class from the ``splitting type'' of $\gamma$.
\end{remark}
\begin{remark}
	Kazhdan--Lusztig conjectured that $\op{KL}$ is injective. This was proven in many cases by Spaltenstein and proven in full by Yun.
\end{remark}
To upgrade the Kazhdan--Lusztig map, we will also need some analog of parabolic subgroups for our situation.
\begin{definition}[parahoric]
	Fix a reductive group $G$ over $\CC$. Then we define the \textit{parahoric subgroups} $P$ as those between $I$ and $LG^\circ$. We then define the \textit{pseudo-Levi subgroup} $L_P\coloneqq P/P^u$, where $P^u$ is the unipotent radical.
\end{definition}
\begin{example}
	One may take $P$ to be $I$ or $L^+G$.
\end{example}
\begin{example}
	For $G=\op{SL}_2$, one can also take $P$ to be the subgroup
	\[\begin{bmatrix}
		t \\ & 1
	\end{bmatrix}L^+\mathrm{SL}_2\begin{bmatrix}
		t^{-1} \\ & 1
	\end{bmatrix}.\]
\end{example}
\begin{remark}
	It turns out that $P/P^u$ is a finite-dimensional reductive group. Its Lie algebra $\mf l_{\mf p}$ sits in a split short exact sequence
	\[0\to\mf p^u\to\mf p\to\mf l_\mf p\to0.\]
\end{remark}
We can now upgrade our Kazhdan--Lusztig map.
\begin{definition}[Kazhdan--Lusztig map]
	Fix a parahoric subgroup $P$. There is a Kazhdan--Lusztig map $\op{KL}_P\colon\mc N_{L_P}\to\op{cl}(W)$ which sends $e\in\mc N_{L_P}$ to the splitting type of some generic $\gamma$ with constant term $\gamma$.
\end{definition}
\begin{remark}
	These Kazhdan--Lusztig maps can be assembled together. Choose parahoric subgroups $P$ and $Q$, and suppose $P\subseteq Q$. Then $P/Q^u$ is a parabolic in $L_Q$, and it admits a natural surjection onto $L_P$. Now, Lusztig--Spaltenstein defined a map $\op{ind}\colon\mc N_{L_P}\to\mc N_{L_Q}$ given by sending $e$ to $e$ plus a generic element in the unipotent radical. It now turns out that
	% https://q.uiver.app/#q=WzAsMyxbMCwwLCJcXG1jIE5fe0xfUH0iXSxbMSwwLCJcXG1jIE5fe0xfUX0iXSxbMSwxLCJcXG9we2NsfShXKSJdLFswLDIsIlxcb3B7S0x9X1AiLDJdLFsxLDIsIlxcb3B7S0x9X1EiXSxbMCwxLCJcXG9we2luZH0iXV0=&macro_url=https%3A%2F%2Fraw.githubusercontent.com%2FdFoiler%2Fnotes%2Fmaster%2Fnir.tex
	\[\begin{tikzcd}[cramped]
		{\mc N_{L_P}} & {\mc N_{L_Q}} \\
		& {\op{cl}(W)}
		\arrow["{\op{ind}}", from=1-1, to=1-2]
		\arrow["{\op{KL}_P}"', from=1-1, to=2-2]
		\arrow["{\op{KL}_Q}", from=1-2, to=2-2]
	\end{tikzcd}\]
	commutes. Thus, there is a map $\colim_P\mc N_{L_P}\to\op{cl}(W)$.
\end{remark}
\begin{example}
	Take $G=\op{GL}_n$, and suppose that $P$ is maximal, and $Q=G$. Then the induction $\op{ind}\colon\mc N_{L_P}\to\mc N_{L_Q}$ basically takes a pair of partitions $(\lambda,\mu)$ to the partition $(\lambda_1+\mu_1,\lambda_2+\mu_2,\ldots)$.
\end{example}
\begin{remark}
	If $P$ and $P'$ are parahoric subgroups with the same Levi (i.e., $L_P=L_{P'}$), then $\mathrm{KL}_{P}=\mathrm{KL}_{{P'}}$. Thus, we may write $\op{KL}_{L_P}$ instead of $\op{KL}_P$.
\end{remark}
Conjecturally, the Kazhdan--Lusztig map knows about strata.
\begin{conj}[Lusztig] \label{conj:lusztig}
	Fix a reductive group $G$ over $\CC$. Then the image of the map $\colim_P\mc N_{L_P}\to\op{cl}(W)$ agrees with the image of the map $\op{Str}(G)\to\op{cl}(W)$ given by choosing the most elliptic map.
\end{conj}
\begin{remark}
	It is also conjectured that the image of $\op{KL}$ agrees with the image of the nilpotent orbits $\mc N$, which is a subset of $\op{Str}(G)$, in $\op{cl}(W)$.
\end{remark}
\begin{remark}
	Roughly speaking, we are hoping to define a map $\colim\mc N_{L_P}\to\op{Str}G$, which we then hope recovers the Kazhdan--Lusztig map once we compose with the map $\op{Str}(G)\to\op{cl}(W)$. Here is a proposal for this map: send $e\in\mc N_{L_P}$ to a stratum containing a generic element of the form $s\exp(e)$, where $s$ is semisimple.
\end{remark}

\subsection{The Minimal Reduction Type}
We now hunt for a map $\op{cl}(W)\to\mc N$, and then we will conjecturally construct a map $\op{cl}(W)\to\op{Str}(G)$.
\begin{definition}[minimal reducton type]
	We define the \textit{minimal reduction type} $\op{RT}_{\min}$ to send a topolotically nilpotent element $\gamma\in\mf g[[t]]$ to the minimal nilpotent orbits $\OO$ of $\mf g$ such that $\gamma$ can be conjugated (by $LG$) into $\OO+t\mf g[[t]]$. We let $\mathrm{RT}_{\min}(\gamma)$ be the collection of these orbits. Then for $[w]\in\op{cl}(W)$, we let $\mathrm{RT}_{\min}([w])$ be $\mathrm{RT}_{\min}(\gamma)$, where $\gamma$ is chosen generically to have splitting type $[w]$.
\end{definition}
\begin{remark}
	In many cases, Yun has shown that $\op{RT}_{\min}(\gamma)$ is a unique nilpotent orbit, and he has computed this in classical types. This is conjectured to be true in all cases. In particular, for generic $\gamma$ with splitting type $[w]$, one finds that $\op{RT}_{\min}(\gamma)$ is a singleton, so $\op{RT}_{\min}\colon\op{cl}(W)\to\mc N$ can be seen to be well-defined.
\end{remark}
\begin{example}
	For $G=\op{GL}_n$, the singleton $\op{RT}_{\min}(\gamma)$ is described using Newton polygons.
\end{example}
\begin{remark}
	These constructions were motivated by affine Springer fibers, which are the finite-dim\-ensional representations $M_\gamma\coloneqq\mathrm H^{\mathrm{top}}(\mathrm{Fl}_\gamma)^{A_\gamma}$. (Lusztig has shown that these admit action from $\widetilde W$.) Now, for $g\in\op{Gr}_\gamma$, one has that $\op{Ad}_{g^{-1}}\gamma\in\mf g[[t]]$ with $\op{Ad}_{g^{-1}}\gamma\equiv e\pmod t$ for some nilpotent $e$, and it turns out that the fiber over $g\in\op{Gr}_\gamma$ in $\op{Fl}_\gamma$ is exactly given by $\mc B_e$. In other words, we have a pullback square as follows.
	% https://q.uiver.app/#q=WzAsNCxbMCwwLCJcXG1jIEJfZSJdLFsxLDAsIlxcb3B7Rmx9X1xcZ2FtbWEiXSxbMCwxLCJcXHtnXFx9Il0sWzEsMSwiXFxvcHtHcn1fe1xcZ2FtbWF9Il0sWzIsMywiIiwwLHsic3R5bGUiOnsidGFpbCI6eyJuYW1lIjoiaG9vayIsInNpZGUiOiJ0b3AifX19XSxbMCwxLCIiLDAseyJzdHlsZSI6eyJ0YWlsIjp7Im5hbWUiOiJob29rIiwic2lkZSI6InRvcCJ9fX1dLFswLDJdLFsxLDNdXQ==&macro_url=https%3A%2F%2Fraw.githubusercontent.com%2FdFoiler%2Fnotes%2Fmaster%2Fnir.tex
	\[\begin{tikzcd}[cramped]
		{\mc B_e} & {\op{Fl}_\gamma} \\
		{\{g\}} & {\op{Gr}_{\gamma}}
		\arrow[hook, from=1-1, to=1-2]
		\arrow[from=1-1, to=2-1]
		\arrow[from=1-2, to=2-2]
		\arrow[hook, from=2-1, to=2-2]
	\end{tikzcd}\]
	Roughly speaking, the uniqueness claim of $\op{RT}_{\min}(\gamma)$ shows that there is a unique Springer representation $E_\OO$ in $M_\gamma$ with longest ``fake degree'' $\dim\mc B_e$.
\end{remark}
The above discussion assembles a map $\op{cl}(W)\to\mc N$.
\begin{notation}
	We define a multivalued map $\op{cl}(W)\to\op{Irr}_*(G)$. We start with $[w]\in\op{cl}(W)$, and then we choose generic topologically nilpotent $\gamma$ of splitting type $[w]$. Then among the $2$-special irreducible representation $E\in\op{Irr}_*(W)$, we let $\widetilde{\mathrm{RT}}_{\min}([w])$ consist of those $E$ appearing in $M_\gamma$ with largest ``fake degree.''
\end{notation}
\begin{conj}[Yun] \label{conj:yun}
	The above construction $\widetilde{\mathrm{RT}}_{\min}([w])$ is a singleton.
\end{conj}
\begin{remark}
	The above conjecture induces a map $\op{cl}(W)\to\op{Irr}_*(W)$, which is in bijection with strata.
\end{remark}
\begin{remark}
	The above conjecture implies \Cref{conj:lusztig}. The main point is that the diagram
	% https://q.uiver.app/#q=WzAsMyxbMCwwLCJcXGNvbGltX1BcXG1jIE5fe0xfUH0iXSxbMSwwLCJcXG9we2NsfShXKSJdLFsxLDEsIlxcb3B7U3RyfShHKSJdLFsxLDIsIlxcd2lkZXRpbGRle1xcbWF0aHJte1JUfX1fe1xcbWlufSJdLFswLDEsIlxcd2lkZXRpbGRle1xcb3B7S0x9fSJdLFswLDIsIihlXFxtYXBzdG8gc1xcZXhwKGUpKSIsMl1d&macro_url=https%3A%2F%2Fraw.githubusercontent.com%2FdFoiler%2Fnotes%2Fmaster%2Fnir.tex
	\[\begin{tikzcd}[cramped]
		{\colim_P\mc N_{L_P}} & {\op{cl}(W)} \\
		& {\op{Str}(G)}
		\arrow["{\widetilde{\op{KL}}}", from=1-1, to=1-2]
		\arrow["{(e\mapsto s\exp(e))}"', from=1-1, to=2-2]
		\arrow["{\widetilde{\mathrm{RT}}_{\min}}", from=1-2, to=2-2]
	\end{tikzcd}\]
	commutes.
\end{remark}
\begin{remark}
	For $\OO\in\mc N$, Chua has shown that $\widetilde{\mathrm{RT}}_{\min}([w])$ is $E_\OO$ if $[w]=\op{KL}(\OO)$. In other words, \Cref{conj:yun} is known if $[w]$ already comes from the Kazhdan--Lusztig map.
\end{remark}
\begin{remark}
	We close the lecture with a question: choose a maximal torus $T\subseteq G_{\CC((t))}$, which we do not assume to be split. Then it is an interesting question to compute $\mathrm{Gr}^{LT^\circ}$. In types $A$ and $C$, these are points; in types $B$ and $D$, these are given by the Lagrangian Grassmannian of certain quadratic spaces. However, this has not been computed for all exceptional types.
\end{remark}

\end{document}